\def\sym#1{\ifmmode^{#1}\else\(\^{#1}\)\fi}
\setlength{\tabcolsep}{5pt}
{\footnotesize
\begin{longtable}{rp{2.2cm}p{2.4cm}p{3.2cm}p{3.8cm}}
\caption*{\textbf{Table AP1a---Authoritarian-Populist Episodes (Strict Subset), N=9}} \\
\hline\hline
No. & Country & Year & Leader & Regime at takeover \\
\hline
\endfirsthead
\hline\hline
No. & Country & Year & Leader & Regime at takeover \\
\hline
\endhead
\hline
\endfoot
1. & Argentina & 1946 & Per\'on (L) & Electoral autocracy (V-Dem=1) \\
2. & Argentina & 1973 & C\'ampora/Per\'on (L) & Electoral autocracy (V-Dem=1) \\
3. & Bolivia & 1952 & Estenssoro (L) & Electoral autocracy (V-Dem=1) \\
4. & Brazil & 1951 & Vargas (L) & Electoral autocracy (V-Dem=1) \\
5. & Chile & 1952 & Ib\'a\~nez (L) & Electoral autocracy (V-Dem=1) \\
6. & Ecuador & 1952 & Velasco Ibarra I (R) & Electoral autocracy (V-Dem=1) \\
7. & Ecuador & 1960 & Velasco Ibarra II (R) & Electoral autocracy (V-Dem=1) \\
8. & Ecuador & 1968 & Velasco Ibarra III (R) & Electoral autocracy (V-Dem=1) \\
9. & Mexico & 1970 & Echeverr\'ia (L) & Electoral autocracy (PRI hegemonic; V-Dem=1) \\
\hline\hline
\multicolumn{5}{p{14cm}}{\textit{Notes:} Strict subset = FST populist episodes in electoral or closed autocracies at the takeover year (V-Dem \texttt{v2x\_regime} $\leq 1$). L = left-wing, R = right-wing (per FST coding). Ecuador comprises three non-consecutive spells of the same leader; see robustness checks.} \\
\end{longtable}
}
